\clearpage
\fancyhead[L]{\small VOYAGE DE DENTRECASTEAUX.}
\hypertarget{native-519-lower}{}\label{native-519-lower}
\begin{center}\small Göttingische gelehrte Anzeigen. Stück 167. pp. 1657..1663. October 21, 1809.\\[3pt]\rule{.72\linewidth}{.3pt}\end{center}

\emph{Voyage de} \textsc{Dentrecasteaux}\emph{, envoyé à la recherche de la Pérouse. Publié par ordre de Sa Majesté l'Empereur et Roi.} \emph{Rédigé par} M. \textsc{de Rossel}\emph{, ancien capitaine de vaisseau.} \emph{Tome second.} \emph{A Paris} 1808. 692 pages in quarto. Paris. The first part of this work, which contains the account of the voyage itself, has already been reviewed in these pages (above, p. 1065). This second volume is devoted entirely to astronomical and nautical observations and is divided into two main parts. The first describes the instruments and examines in general how much accuracy can be achieved with them and under what circumstances observations can best be made. The second presents in the greatest detail, arranged in tables, the astronomical observations made during the voyage on the frigate La Recherche, including some made on the other frigate, L'Esperance. The astronomical observations that matter so much to the safety of navigation, and the work of determining latitude and longitude from them, usually fall to people who have had no mathematical training as such and have generally acquired little more than a mechanical skill in observing and calculating. That is the context of the treatise that forms the first part. It is certainly a very valuable work by an observer as knowledgeable and experienced as Captain \textsc{Rossel}, and aspiring seamen would do well to study it to sharpen their judgment of the value to assign to 'observations. Amateurs who observe on land will also read it with profit. One can only wish that this treatise did not form part of so expensive a work, with which it really has no direct connection, or that it were at least made available to more readers in a separate printing. We would also have liked smaller subdivisions to make the work easier to consult. There are sections, to be sure, but sections 119 pages long. Of the two chapters into which the treatise is divided, the first deals with the instruments carried by the two frigates and also explains the particular merits of each instrument. Besides several \textsc{Bordai} reflecting circles, each frigate carried an astronomical repeating circle for use on land, a marine chronometer, an astronomical clock, a counter, an achromatic telescope, a dip \SourceLeaf{526}compass of a new design, and an azimuth compass. At several places the dip compasses were also used to test the duration of the needle's oscillations in the magnetic meridian, and \textsc{Rossel} includes a table for reducing the oscillations to infinitesimally small ones. The whole apparatus was enclosed in a case with glass panes, but the case had to be opened for these experiments. \textsc{Rossel} therefore thinks that the movement of the air may sometimes have made the experiments somewhat uncertain. Finally, the author describes the construction of an astronomical tent of the kind carried by each ship for use in observations on land. — The second chapter explains the nature and limits of errors arising from the instruments, the observations themselves, or the tables, and how to correct them or at least reduce them as far as possible. The various kinds of observations are discussed here in 9 sections. §. I. Determining the time from individual altitudes. Among other things, the author draws attention to the effect of errors made in determining the dip of the horizon. He therefore recommends always observing from as high above the surface of the sea as practicable. The same remark clearly applies to all angles of altitude observed at sea. §. II. Determining latitude from meridian altitudes. \textsc{Rossel} considers solar altitudes near the zenith accurate to within only 2..3 minutes. His reason for regarding the measurement as less reliable in these circumstances does not quite convince us. He says that when the reflecting circle is rocked, the image of the Sun separates only very slowly from the sea horizon, so that the actual vertical circle of contact cannot be determined accurately. Attempts to take stellar altitudes at sea never succeeded. Even under the most favorable conditions, errors of up to 5 minutes could not be avoided. This may be due in part to the imperfect optics of the French instruments, whose telescopes have much too small an aperture. \textsc{Rossel} puts the accuracy of latitudes determined from lunar altitudes at 3 minutes by day and 4 minutes by night. But this includes 1 minute for error in the Moon's declination taken from the ephemerides. That error could easily be avoided by taking the trouble to calculate the declination from the longitude and latitude, which the ephemerides usually give to the second. §. III. Determining latitude from two solar altitudes off the meridian. The indirect methods of solution are rightly recommended in preference to the direct one, even though the direct solution does not necessarily require calculations as lengthy and laborious, or so complicated a distinction among many cases, as \textsc{Rossel} believes. There is nothing new here about the solution itself. But the effects of errors in the observations, in estimating how far the ship has moved between the observations, and in the rate of the clock are discussed in great detail. The final result is that with this method one can always determine latitude to within 4 to 5 minutes, even in the worst case, when all the errors conspire in the same direction, provided the necessary precautions are taken in selecting the observations. What is said about the relative accuracy of the various indirect methods of calculation may be rather superfluous in the cases that arise in practice. §. IV. Observing the Sun's azimuth to determine the declination of the magnetic needle. \textsc{Rossel} points out that the declination measured aboard ship off the coast sometimes differs considerably from the value found a short distance away on land. The observations near Tenerife that we give below offer a striking example. §. V. Determining longitude from measured distances of the Moon from the Sun and from fixed stars. This section contains several valuable practical remarks that carry all the more weight because they come from long experience. The author is surely right in saying, “Anything that can make so delicate an observation easier is valuable. We know of \SourceLeaf{527}none that calls for greater natural aptitude. We are therefore firmly convinced that when scholars have done everything their genius can suggest to perfect this method, the observer’s skill will still have much to add.” — In calculating lunar parallax to reduce the measured distances, \textsc{Rossel} takes the spheroidal shape of the Earth into account only insofar as the horizontal parallax at the latitude of the place of observation is smaller than the equatorial parallax. He neglects the inclination of the vertical to the radius drawn toward the Earth's center. Admittedly, this makes little difference for observations of the kind that can be made with ordinary French reflecting circles. But in our opinion \textsc{Rossel} has misjudged the effect of this last circumstance. He considers only that the parallax is calculated for an incorrect \emph{altitude}, and forgets that its \emph{direction} is also affected. We therefore think he is mistaken in concluding that this somewhat inaccurate method of calculation distorts the result most when the Moon is near culmination. As a rule, the error is greatest precisely when the Moon is near the prime vertical. §. VI. Determining longitude by marine chronometers. §. VII. How the longitudes determined by lunar distances must be combined with the differences of longitude obtained from marine chronometers. §. VIII. Determining longitude from occultations of stars and eclipses of Jupiter's satellites. §. IX. Determining the azimuths of terrestrial objects (relevemens).

The second part of the work contains the observation journal itself. It is preceded by a table of the errors of \textsc{Mason}'s lunar tables based on simultaneous Greenwich observations, and by the corrections derived from it for the longitudes determined during the voyage. To single out any of the observational results that are so important for the geography of the South Seas would be outside the plan of these \emph{Göttingische gelehrte Anzeigen}. It would be all the less necessary, since a very full extract of the geographical determinations is already available in the \emph{Monatliche Correspondenz}. Still, some readers may be glad to see here the few determinations made with the magnetic needle that were omitted in that journal. At \emph{Brest}, on September 20, 1791, the northward dip of the needle was found to be $71^\circ\,30'$ and the duration of an infinitesimally small oscillation $2^{\mathrm s}02$. At \emph{Santa Cruz} on Tenerife, on October 21, 1791, the northward dip was $62^\circ\,25'$ and the mean period of oscillation $2^{\mathrm s}081$. The declination of the needle aboard the Recherche was $18^\circ\,7'$ N. W., at the observatory $21^\circ\,33'$, and on the mole $23^\circ\,43'$. At \emph{Port du Nord} in Van Diemen's Land, on May 11, 1792, the dip was $70^\circ\,50'$ southward, the mean period of oscillation $1^{\mathrm s}869$, and the declination $8^\circ\,1'\,12''$ N. E. On the island of \emph{Amboina}, on October 9, 1792, the dip was $20^\circ\,37'$ southward and the period of oscillation $2^{\mathrm s}403$. At \emph{Port au Sud} in Van Diemen's Land, on February 7, 1793, the declination was $2^\circ\,33'\,47''$ N. E. with one needle and $3^\circ\,10'\,24''$ with the other. The dips were $72^\circ\,22'$ and $70^\circ\,48'$ southward, and the periods of oscillation $1^{\mathrm s}8498$ and $1^{\mathrm s}817$. At \emph{Sourabaya} on the island of Java, on May 9, 1794, the dip was $25^\circ\,40'$ southward, and the period of oscillation was $2^{\mathrm s}429$ with one needle and $2^{\mathrm s}485$ with the other. The various observations of the duration of the oscillations thus seem to show that the intensity of the magnetic force at the equator is considerably weaker than at higher latitudes.
